% 原书物理页=印刷页38--44；1.12 全节.（1.12.10 推论及其证明实际印于原书第 45 页顶部、1.13 节标题之前，属本节内容，一并译出附于末尾.）
\section{有限维表示}\label{sec:lie-1-12}

\begin{numpara}[表示的范畴]\label{np:lie-1-12-1}
记 $\mathbb K$ 为域 $\mathbb R$ 或 $\mathbb C$，并回顾以下基本定义.
\begin{enumerate}[label=\roman*)]
\item 设 $G$ 是一个 Lie 群.
\begin{itemize}[label=--]
\item 一个\emph{以 $\mathbb K$ 为系数的表示}是一个偶对 $(V,\rho)$，其中 $V$ 是一个 $\mathbb K$-向量空间，而 $\rho:G\longrightarrow\mathrm{GL}(V)$ 是一个光滑态射\footnote{或连续；前面已经看到二者是等价的.}. 我们常简化记号，对 $g\in G$ 与 $v\in V$ 令 $gv=\rho(g)(v)$.
\item $(V,\rho)$ 与 $(V',\rho')$ 之间的一个\newterm{lie-morphism-of-representations}{表示态射}{morphism of representations}（以 $\mathbb K$ 为系数）是一个 $\mathbb K$-线性映射 $f:V\longrightarrow V'$，使得 $\forall g\in G$，$\rho'(g)\circ f=f\circ\rho(g)$；用简化记号即 $\forall g\in G$，$\forall v\in V$，$f(gv)=g\,f(v)$. 我们把由这些态射构成的 $\mathbb K$-向量空间记作 $\operatorname{Hom}_{\mathbb K G}(V,V')$，在没有歧义时可以省略 $\mathbb K$.
\end{itemize}
\item 设 $\mathfrak g$ 是一个 $\mathbb R$-Lie 代数.
\begin{itemize}[label=--]
\item 一个\emph{以 $\mathbb K$ 为系数的表示}是一个偶对 $(V,r)$，其中 $V$ 是一个 $\mathbb K$-向量空间，而 $r:\mathfrak g\longrightarrow\mathrm{GL}(V)$ 是一个 $\mathbb R$-Lie 代数态射. 我们常简化记号，对 $X\in\mathfrak g$ 与 $v\in V$ 令 $Xv=r(X)(v)$.
\item $(V,r)$ 与 $(V',r')$ 之间的一个以 $\mathbb K$ 为系数的表示态射是一个 $\mathbb K$-线性映射 $f:V\longrightarrow V'$，使得 $\forall X\in\mathfrak g$，$r'(X)\circ f=f\circ r(X)$；用简化记号即 $\forall X\in\mathfrak g$，$\forall v\in V$，$f(Xv)=X\,f(v)$. 我们把由这些态射构成的 $\mathbb K$-向量空间记作 $\operatorname{Hom}_{\mathbb K\mathfrak g}(V,V')$，在没有歧义时可以省略 $\mathbb K$.
\end{itemize}
\end{enumerate}
\end{numpara}
\begin{remark*}
若干自然的例子取实系数（例如伴随表示），但我们将会看到，取复系数时理论更为宜人.
\end{remark*}
有了这些定义，我们得到两个范畴 $\operatorname{Rep}_{\mathbb K}(G)$ 与 $\operatorname{Rep}_{\mathbb K}(\mathfrak g)$，其对象为有限维表示. 易见这两个范畴都是\newterm{lie-abelian-category}{阿贝尔范畴}{abelian category}，而且核、余核与直和同通向 $\mathbb K$-向量空间范畴的\newterm{lie-forgetful-functor}{遗忘函子}{forgetful functor}交换. 换言之，若 $f:V\longrightarrow V'$ 是表示之间的一个态射，那么它的核是 $V$ 的一个\newterm{lie-subrepresentation}{子表示}{subrepresentation}，它的像是 $V'$ 的一个子表示，而它的余核带有唯一的表示结构，使得从 $V'$ 出发的投影连续（``\newterm{lie-quotient-representation}{商表示}{quotient representation}''）. 同样，我们把两个表示的直和定义为底空间的直和，并配备 $G$ 或 $\mathfrak g$ 的显然作用.
\begin{exercise*}
验证上一句话中的各个论断.
\end{exercise*}

\begin{numpara}[``求导''函子]\label{np:lie-1-12-2}
若 $(V,\rho)$ 是 $G$ 的一个表示，我们可以通过求导给它配上 $\operatorname{Lie}(G)$ 的一个表示 $(V,d\rho)$. 此外，若 $f:V\longrightarrow V'$ 是 $G$ 的表示之间的一个态射，那么它也是相应的 $\operatorname{Lie}(G)$ 表示之间的一个态射，这是因为在 $g=e$ 处求等式 $f\circ\rho(g)=\rho'(g)\circ f$ 的微分便得到 $\forall X\in\operatorname{Lie}(G)$，$f\circ d\rho(X)=d\rho'(X)\circ f$. 这样我们便得到一个正合函子
\[
\begin{array}{rcl}
\operatorname{Rep}_{\mathbb K}(G) & \longrightarrow & \operatorname{Rep}_{\mathbb K}(\operatorname{Lie}(G))\\
(V,\rho) & \longmapsto & (V,d\rho)
\end{array},
\]
它是研究 $G$ 的表示的最重要的工具.
\begin{proposition*}
上面的``求导''函子：
\begin{enumerate}[label=\roman*)]
\item 当 $G$ 连通时是\newterm{lie-fully-faithful}{全忠实的}{fully faithful}；
\item 当 $G$ 单连通时是\newterm{lie-essentially-surjective}{本质满的}{essentially surjective}（从而是一个范畴等价）.
\end{enumerate}
\end{proposition*}
\begin{proof}
i) 要证明的是：对 $\operatorname{Rep}_{\mathbb K}(G)$ 中所有的 $V,V'$，映射 $\operatorname{Hom}_{G}(V,V')\longrightarrow\operatorname{Hom}_{\operatorname{Lie}(G)}(V,V')$ 都是双射. 单性是重言式的. 为证满性，我们利用公式 $\rho(\exp_{G}(X))=\exp_{\mathrm{GL}(V)}(d\rho(X))$，它表明：若 $f\in\operatorname{Hom}_{\operatorname{Lie}(G)}(V,V')$，则对 $\exp_{G}(\operatorname{Lie}(G))$ 中的一切 $g$ 都有 $f\circ\rho(g)=\rho'(g)\circ f$. 而使这个等式成立的 $g\in G$ 全体构成一个子群. 但由于 $G$ 连通，它由 $\exp_{G}(\operatorname{Lie}(G))$ 生成，故该等式对一切 $g$ 成立.
ii) 把定理~\ref{thm:lie-1-9-2} 之 ii) 应用于 $r:\operatorname{Lie}(G)\longrightarrow\mathfrak{gl}(V)$ 即得.
\end{proof}
\begin{remark*}
在这里局限于有限维表示是重要的. $\mathfrak{sl}_2(\mathbb C)$ 存在不能积分成 $\mathrm{SL}_2(\mathbb C)$ 的无限维表示.
\end{remark*}
\end{numpara}

\begin{numpara}[复变体与 Weyl 酉技巧]\label{np:lie-1-12-3}
本小节中 $\mathbb K=\mathbb C$.
\begin{enumerate}[label=\roman*)]
\item 若 $G$ 是一个复 Lie 群，则 $G$ 的一个\newterm{lie-holomorphic-representation}{全纯表示}{holomorphic representation}是一个复表示 $(V,\rho)$，其中 $\rho$ 全纯. 这些表示构成 $\operatorname{Rep}_{\mathbb C}(G)$ 的一个\newterm{lie-full-subcategory}{全子范畴}{full subcategory} $\operatorname{Rep}^{\mathrm{hol}}_{\mathbb C}(G)$.
\item 若 $\mathfrak g$ 是一个 $\mathbb C$-Lie 代数，我们依类比把 $\mathfrak g$ 的一个复表示 $(V,r)$ 称为``全纯的''\footnote{这一术语并非惯用，但在这里是相当合理的. 我们稍后会对它稍作修改\ldots}，如果 $r$ 是 $\mathbb C$-线性的. 同样得到 $\operatorname{Rep}_{\mathbb C}(\mathfrak g)$ 的一个全子范畴 $\operatorname{Rep}^{\mathrm{hol}}_{\mathbb C}(\mathfrak g)$.
\item 求导函子把 $\operatorname{Rep}^{\mathrm{hol}}_{\mathbb C}(G)$ 送入 $\operatorname{Rep}^{\mathrm{hol}}_{\mathbb C}(\operatorname{Lie}(G))$，当 $G$ 连通时全忠实，当 $G$ 单连通时本质满.
\end{enumerate}
\emph{复化.}\ 从一个 $\mathbb R$-Lie 代数 $\mathfrak g$ 出发. 任何复表示 $r:\mathfrak g\longrightarrow\mathfrak{gl}(V)$ 都唯一地延拓为一个``全纯''表示 $r_{\mathbb C}:\mathfrak g_{\mathbb C}\longrightarrow\mathfrak{gl}(V)$. 事实上，数量扩张的泛性质表明，如此得到的函子
\[
\operatorname{Rep}_{\mathbb C}(\mathfrak g)\longrightarrow\operatorname{Rep}^{\mathrm{hol}}_{\mathbb C}(\mathfrak g_{\mathbb C})
\]
是一个范畴等价，它是限制函子
\[
\operatorname{Rep}^{\mathrm{hol}}_{\mathbb C}(\mathfrak g_{\mathbb C})\longrightarrow\operatorname{Rep}_{\mathbb C}(\mathfrak g).
\]
的逆.
同样，若从 Lie 群 $G$ 及其一个复化 $G\xrightarrow{\ \iota\ }G_{c}$ 出发，则通过限制得到一个函子
\[
\operatorname{Rep}^{\mathrm{hol}}_{\mathbb C}(G_{c})\xrightarrow{\ \mathrm{res}\ }\operatorname{Rep}_{\mathbb C}(G).
\]
复化的泛性质表明该函子本质满. 过渡到 Lie 代数上可以看出，它还是全忠实的，因而当这些群连通时是一个范畴等价.
\begin{example*}
（Weyl 酉技巧）看下图
\[
\begin{tikzcd}
\operatorname{Rep}_{\mathbb C}(\mathrm{SL}_n(\mathbb R)) &
\operatorname{Rep}^{\mathrm{hol}}_{\mathbb C}(\mathrm{SL}_n(\mathbb C))
\ar[l,"\mathrm{res}"'] \ar[r,"\mathrm{res}"] &
\operatorname{Rep}_{\mathbb C}(\mathrm{SU}(n)) \\
\operatorname{Rep}_{\mathbb C}(\mathfrak{sl}_n(\mathbb R)) &
\operatorname{Rep}^{\mathrm{hol}}_{\mathbb C}(\mathfrak{sl}_n(\mathbb C))
\ar[l,"\mathrm{res}"'] \ar[r,"\mathrm{res}"] &
\operatorname{Rep}_{\mathbb C}(\mathfrak{su}(n))
\ar[from=1-1,to=2-1]
\ar[from=1-2,to=2-2,"\simeq"]
\ar[from=1-3,to=2-3,"\simeq"]
\end{tikzcd}
\]
我们刚刚看到，所有限制函子都是范畴等价. $\mathrm{SU}(n)$ 与 $\mathrm{SL}_n(\mathbb C)$ 的两个求导函子也是范畴等价，因为这两个群单连通. 由此推出 $\mathrm{SL}_n(\mathbb R)$ 的求导函子也是范畴等价，尽管后者不是单连通的. 这样我们就把 $\mathrm{SL}_n(\mathbb R)$ 的表示的研究化归为群 $\mathrm{SU}(n)$ 的表示的研究，而后者是\emph{紧}的.
\end{example*}
由此还可推出如下性质：$\mathrm{SL}_n(\mathbb R)$ 的任何非平凡覆盖都不能嵌入任何 $\mathrm{GL}_n(\mathbb C)$ 之中.
\end{numpara}

\begin{numpara}[平凡表示、单位表示与不变向量]\label{np:lie-1-12-4}
设 $V$ 是 Lie 群 $G$ 或 Lie 代数 $\mathfrak g$ 的一个表示的底向量空间. 我们记
\begin{itemize}[label=--]
\item （群的情形）：$V^{G}:=\{v\in V,\ \forall g\in G,\ gv=v\}$，即 $G$ 作用下不变的\newterm{lie-invariant-vector}{不变向量}{invariant vector}所成的向量子空间.
\item （代数的情形）：$V^{\mathfrak g}:=\{v\in V,\ \forall X\in\mathfrak g,\ Xv=0\}$，即被 $\mathfrak g$ 零化的向量所成的向量子空间.
\end{itemize}
\begin{exercise*}
设 $(V,\rho)$ 是\emph{连通}群 $G$ 的表示，证明 $V^{\operatorname{Lie}(G)}=V^{G}$.
\end{exercise*}
$G$ 的表示 $V$（相应地，$\mathfrak g$ 的表示 $V$）称为\newterm{lie-trivial-representation}{平凡的}{trivial}，如果 $V=V^{G}$（相应地，$V=V^{\mathfrak g}$）.
$G$（相应地，$\mathfrak g$）的\newterm{lie-unit-representation}{单位表示}{unit representation}是配备平凡表示的向量空间 $\mathbb K$；在没有歧义时我们简单地把它记作 $\mathbb K$.
\end{numpara}

\begin{numpara}[张量积、逆步表示与内 Hom]\label{np:lie-1-12-5}
设 $V,V'$ 是两个 $\mathbb K$-向量空间，分别是 Lie 群 $G$ 的表示 $\rho,\rho'$ 或 Lie 代数 $\mathfrak g$ 的表示 $r,r'$ 的底空间.
\begin{enumerate}[label=\roman*)]
\item 我们在 $V\otimes_{\mathbb K}V'$ 上配备如下定义的\newterm{lie-tensor-product-representation}{张量积表示}{tensor product representation}：
\begin{itemize}[label=--]
\item （群的情形）$g\mapsto\rho(g)\otimes\rho'(g)$，或者写成 $g(v\otimes v'):=gv\otimes gv'$，$\forall v\in V$，$v'\in V'$；
\item （代数的情形）$X\mapsto r(X)\otimes\mathrm{id}_{V'}+\mathrm{id}_{V}\otimes r'(X)$，或者写成 $X(v\otimes v'):=Xv\otimes v'+v\otimes Xv'$，$\forall v\in V$，$v'\in V'$.
\end{itemize}
\begin{exercise*}
验证第二个定义确实给出 $\mathfrak g$ 的一个表示；再验证 $\rho\otimes\rho'$ 的导数确为 $d\rho\otimes\mathrm{id}+\mathrm{id}\otimes d\rho'$.
\end{exercise*}
\begin{exercise*}
若 $V'=\mathbb K$ 是单位表示，验证 $V\otimes_{\mathbb K}\mathbb K\simeq V$（这也解释了``单位表示''这一术语）.
\end{exercise*}
\item 我们在 $\operatorname{Hom}_{\mathbb K}(V,V')$ 上配备如下定义的\newterm{lie-internal-hom}{内 Hom 表示}{internal Hom}：
\begin{itemize}[label=--]
\item （群的情形）$(g,f)\mapsto\rho'(g)\circ f\circ\rho(g^{-1})$，或者写成 $(gf)(v):=g\,f(g^{-1}v)$，$\forall v\in V$；
\item （代数的情形）$(X,f)\mapsto r'(X)\circ f-f\circ r(X)$，或者写成 $(Xf)(v):=Xf(v)-f(Xv)$，$\forall v\in V$.
\end{itemize}
我们把这一 Hom-内表示记作 $\mathcal{H}om_{\mathbb K}(V,V')$（用花体的 $\mathcal H$）.
\begin{exercise*}
验证这些公式确实定义了表示，并且 $G$ 的内 Hom 的导数就是 $\operatorname{Lie}(G)$ 的内 Hom.
\end{exercise*}
\begin{remark*}
我们有 $\operatorname{Hom}_{G}(V,V')=\mathcal{H}om_{\mathbb K}(V,V')^{G}$ 与 $\operatorname{Hom}_{\mathfrak g}(V,V')=\mathcal{H}om_{\mathbb K}(V,V')^{\mathfrak g}$.
\end{remark*}
\item ii) 的一个重要特殊情形是 $V'=\mathbb K$ 为单位表示的情形. 此时我们简单地记 $V^{\vee}:=\mathcal{H}om_{\mathbb K}(V,\mathbb K)$，并称之为 $V$ 的\newterm{lie-contragredient}{逆步表示}{contragredient representation}（或称对偶表示）.
\end{enumerate}
张量积与内 Hom 之间有如下联系：
\begin{lemma*}
态射 $V^{\vee}\otimes_{\mathbb K}V'\longrightarrow\mathcal{H}om_{\mathbb K}(V,V')$，它把 $v^{*}\otimes v'$ 送到态射 $(v,v')\mapsto\langle v,v^{*}\rangle v'$，是表示之间的同构.
\end{lemma*}
\begin{proof}
习题.
\end{proof}
\end{numpara}

\begin{numpara}[表示的约化]\label{np:lie-1-12-6}
一个表示 $V$（Lie 群的或 Lie 代数的）称为：
\begin{itemize}[label=--]
\item \newterm{lie-irreducible-representation}{不可约的}{irreducible}，如果它仅有的子表示是 $\{0\}$ 与 $V$；
\item \newterm{lie-completely-reducible}{完全可约的}{completely reducible}，如果它是不可约子表示的直和；
\item \newterm{lie-indecomposable}{不可分解的}{indecomposable}，如果它不是两个子表示的直和.
\end{itemize}
\begin{example*}
（$\mathfrak g=\mathbb R$ 的情形）$\mathfrak g=\mathbb R$ 的以 $\mathbb K$ 为系数的表示 $(V,r)$ 由自同态 $X:=r(1)\in\operatorname{End}_{\mathbb K}(V)$ 决定. 当 $\mathbb K=\mathbb C$ 时可以看到：
\begin{itemize}[label=--]
\item $V$ 不可约当且仅当 $\dim_{\mathbb C}(V)=1$（本征直线的存在性.）
\item $V$ 完全可约当且仅当 $X$ 可对角化.
\item $V$ 不可分解当且仅当在 Jordan 分解 $X=X_{s}+X_{n}$ 中，$X_{n}$ 的幂零阶数等于 $V$ 的维数（此时 $X_{s}$ 是一个位似）.
\end{itemize}
\emph{相反，}当 $\mathbb K=\mathbb R$ 时，$\mathfrak g=\mathbb R$ 的不可约表示可以是二维的. 例如由
\[
r(1)=\begin{pmatrix}0&1\\-1&0\end{pmatrix}
\]
给出的表示便是如此.
\end{example*}
\begin{exercise*}
设 $(V,\rho)$ 是连通 Lie 群 $G$ 的一个有限维表示. 证明：$(V,\rho)$ 不可约（相应地，不可分解、完全可约）当且仅当 $(V,d\rho)$ 如此.
\end{exercise*}
\begin{exercise*}
证明 $V$ 不可约、不可分解或完全可约当且仅当 $V^{\vee}$ 如此.
\end{exercise*}
容易用对维数的递降归纳法看出，任何有限维表示都是不可分解子表示的直和. 于是，要理解范畴 $\operatorname{Rep}_{\mathbb K}(G)$ 或 $\operatorname{Rep}_{\mathbb K}(\mathfrak g)$，先验地需要对不可分解表示进行分类. 在上面的例子中，它们由维数与 $X$ 的唯一特征值分类. 但一般说来，这个问题可能变得过于困难，甚至毫无希望，例如对阿贝尔 Lie 代数 $\mathfrak g=\mathbb R^{3}$（此时一个表示相当于在某个空间 $V$ 上给定两两交换的 3 个自同态）.
如果我们知道 $G$ 或 $\mathfrak g$ 的表示全都\emph{完全可约}，情形就会大为简化. 此时只需对不可约表示进行分类，而这往往更容易.
我们将会看到，这一完全可约性对``半单的''经典群是满足的. 在此之前，我们先证明它对\emph{紧}群成立.
\begin{lemma*}
对一个表示 $V$，下列性质等价：
\begin{enumerate}[label=\roman*)]
\item $V$ 的任何子表示 $W$ 都容许一个稳定补（从而是一个子表示）$W'$，使得 $W'\oplus W=V$；
\item $V$ 完全可约；
\item $V$ 是它的不可约子表示之和.
\end{enumerate}
\end{lemma*}
\begin{proof}
i) $\Rightarrow$ ii). 对 $V$ 的维数作归纳. 若 $V$ 不可约（因而特别地当 $\dim(V)=1$ 时）则无事可做. 否则，取 $V$ 的一个维数 $>0$ 的极小子表示 $W$. 它于是不可约，且由假设容许一个稳定补 $V'$. 我们来证明 $V'$ 仍满足 i). 设 $W'\subset V'$ 是一个子表示，并设 $W''$ 是 $W\oplus W'$ 在 $V$ 中的一个稳定补. 那么它沿 $W$ 到 $V'$ 上的投影是 $W'$ 在 $V'$ 中的一个稳定补. 这样 $V'$ 满足 i)，由归纳假设，$V$ 是不可约表示的直和. 故 $V=W\oplus V'$ 也成立.
ii) $\Rightarrow$ iii) 是重言式的.
iii) $\Rightarrow$ i). 我们对 $W$ 的余维数作归纳，来证明 $W\subset V$ 容许一个稳定补. 若余维数为零则无事可做. 否则，由假设存在 $V$ 的一个不可约表示 $U$ 使得 $U\nsubseteq W$，从而 $W\cap U=\{0\}$. 把归纳假设用于 $W\oplus U$，便得到 $W\oplus U$ 的一个补 $W'_{1}$，于是只需令 $W':=U\oplus W'_{1}$.
\end{proof}
\begin{lemma*}
设 $V$ 是一个复表示（相应地，实表示），其上配备一个在下述意义下不变的埃尔米特内积（相应地，欧氏内积）：
\begin{itemize}[label=--]
\item （群的情形）$\forall v,v'\in V$，$\forall g\in G$，$\langle gv,gv'\rangle=\langle v,v'\rangle$；
\item （代数的情形）$\forall v,v'\in V$，$\forall X\in\mathfrak g$，$\langle Xv,v'\rangle+\langle v,Xv'\rangle=0$.
\end{itemize}
那么 $V$ 完全可约.
\end{lemma*}
\begin{proof}
事实上，设 $W\subset V$ 是一个稳定子空间. 那么它的正交补 $W^{\perp}=\{v\in V,\ \forall w\in W,\ \langle v,w\rangle=0\}$ 与 $W$ 构成直和 $W\oplus W^{\perp}$（因为 $\langle\ ,\ \rangle$ 正定），并且是稳定的.
\end{proof}
\begin{remark*}
若记 $\Psi(v,w)=\langle v,w\rangle$，则``不变''这一条件在复情形等价于 $G\subset\mathrm{U}(\Psi)$（群的情形）或 $\mathfrak g\subset\mathfrak u(\Psi)$（代数的情形），在实情形等价于 $G\subset\mathrm{O}(\Psi)$ 或 $\mathfrak g\subset\mathfrak{so}(\Psi)$.
\end{remark*}
\end{numpara}
\begin{corollary}\label{cor:lie-1-12-7}
设 $G$ 是一个紧 Lie 群. 那么 $G$ 的所有有限维表示都完全可约.
\end{corollary}
\begin{proof}
作为 Haar 测度存在性的应用，我们已经看到：对 $\mathrm{GL}(V)$ 的任何紧子群 $H$，$V$ 上都存在一个 $H$ 不变的埃尔米特/欧氏内积. 只需把它用于所考虑的表示下 $G$ 在 $\mathrm{GL}(V)$ 中的像，再应用前一个引理即可.
\end{proof}
\begin{example*}
（$G=S^{1}$ 的情形）$S^{1}$ 的一个复不可约表示通过求导给出 $\operatorname{Lie}(G)=\mathbb R$ 的一个复不可约表示，因而是一维的，由一个同态 $S^{1}\longrightarrow\mathrm{GL}_1(\mathbb C)=\mathbb C^{\times}$ 给出. 若把 $S^{1}$ 等同于 $\{z\in\mathbb C,\ |z|=1\}$，这样的同态形如 $z\mapsto z^{n}$，其中 $n\in\mathbb Z$. 于是，$S^{1}$ 的任何表示 $V$ 都同构于一个直和 $\bigoplus_{n\in\mathbb Z}V_{n}^{\oplus e_{n}}$，其中 $V_{n}=(\mathbb C,\ z\mapsto z^{n})$，而 $e_{n}$ 是 $V_{n}$ 作为 $V$ 的子表示出现时的``重数''.
\end{example*}
\emph{应用.}\ $\mathrm{SU}(n)$ 的所有复表示都完全可约. 因而 $\mathfrak{su}(n)$、$\mathfrak{sl}_n(\mathbb R)$ 与 $\mathrm{SL}_n(\mathbb R)$ 的表示，以及 $\mathrm{SL}_n(\mathbb C)$ 与 $\mathfrak{sl}_n(\mathbb C)$ 的全纯表示也是如此. 只需对其中一个对象的不可约表示进行分类，便能分类其余对象的不可约表示. 我们将在下一小节对 $\mathfrak{sl}_2(\mathbb R)$ 做这件事.

\begin{numpara}[Schur 引理]\label{np:lie-1-12-8}
设 $f:V\longrightarrow V'$ 是不可约表示之间的一个态射. 若 $f$ 非零，则 $\operatorname{Ker}(f)\neq V$，故由 $V$ 的不可约性有 $\operatorname{Ker}(f)=0$；同样 $\operatorname{Im}(f)\neq 0$，故由 $V'$ 的不可约性有 $\operatorname{Im}(f)=V'$. 于是 $f$ 或为零，或可逆.
特别地，若 $V$ 是有限维不可约表示，则 $\operatorname{End}_{\mathbb K G}(V)$ 是一个有限维 $\mathbb K$-可除代数. 我们知道，当 $\mathbb K=\mathbb C$ 时，这蕴含 $\operatorname{End}_{\mathbb C G}(V)=\mathbb C$（Schur 引理）. 当 $\mathbb K=\mathbb R$ 时，可能有 $\operatorname{End}_{\mathbb R G}(V)=\mathbb R$、$\mathbb C$ 或 $\mathbb H$.
\emph{应用.}\ 一个完全可约表示 $V$ 可以\emph{以唯一的方式}（在相差次序的意义下）写成 $V=V_{1}^{\oplus e_{1}}\oplus\cdots\oplus V_{r}^{\oplus e_{r}}$，其中 $V_{1},\cdots,V_{r}$ 是两两不同构的不可约表示. 事实上，重数 $e_{i}$ 可作为 $\dim_{\mathbb K}\bigl(\operatorname{Hom}_{\mathbb K G}(V_{i},V)\bigr)/\dim_{\mathbb K}\bigl(\operatorname{End}_{\mathbb K G}(V_{i})\bigr)$ 恢复出来.
\end{numpara}
\ref{sec:lie-1-12} 的末尾部分未在课上讲授.

\begin{numpara}[乘积的表示]\label{np:lie-1-12-9}
设 $G,G'$ 是两个群，或者 $\mathfrak g,\mathfrak g'$ 是两个 Lie 代数. 若 $W$ 是 $G$（相应地，$\mathfrak g$）的一个复表示，而 $W'$ 是 $G'$（相应地，$\mathfrak g'$）的一个表示，我们在 $W\otimes_{\mathbb C}W'$ 上配备表示 $(g,g')(v\otimes v'):=(gv,gv')$\footnote{译注：按上文张量积表示的定义，此处右端应为 $gv\otimes g'v'$；原书印作 $(gv,gv')$，照录.}（$G\times G'$ 的情形），或相应地 $(X,X')(v\otimes v'):=Xv\otimes v'+v\otimes Xv'$（$\mathfrak g\oplus\mathfrak g'$ 的情形）.
\begin{proposition*}
在上述两种情形（群或代数）中，都有：
\begin{enumerate}[label=\roman*)]
\item 若 $W$ 与 $W'$ 不可约，则 $W\otimes_{\mathbb C}W'$ 也不可约（\textbf{就乘积而言！}）；
\item $G\times G'$（相应地，$\mathfrak g\oplus\mathfrak g'$）的任何复不可约表示 $V$ 都具有上文形状 $W\otimes_{\mathbb C}W'$；
\item 若 $W$ 与 $W'$ 完全可约，则 $W\otimes_{\mathbb C}W'$ 也完全可约.
\end{enumerate}
\end{proposition*}
\begin{proof}
为了固定记号，我们处理群的情形；代数的情形完全类似. 先做一个一般性的构造. 若 $W\in\operatorname{Rep}_{\mathbb C}(G)$ 且 $V\in\operatorname{Rep}_{\mathbb C}(G\times G')$，我们在 $\operatorname{Hom}_{\mathbb C G}(W,V|_{G})$ 上配备由 $(g'f)(w):=(e,g')f(w)$ 定义的 $G'$-表示（这是``内 Hom''构造的一个变体）. 于是有一个求值映射：
\[
\operatorname{ev}_{W,V}:\quad
\begin{array}{rcl}
W\otimes_{\mathbb C}\operatorname{Hom}_{\mathbb C G}(W,V) & \longrightarrow & V\\
w\otimes f & \longmapsto & f(w)
\end{array}.
\]
容易验证它是 $G\times G'$ 的表示之间的态射（习题：请验证！）.
现在设 $V$ 与 $W$ 不可约，且 $W$ 同构于 $V|_{G}$ 的一个不可约子表示. 此时 $\operatorname{ev}_{W,V}$ 非零，从而是满射. 令 $n':=\dim_{\mathbb C}\bigl(\operatorname{Hom}_{\mathbb C G}(W,V|_{G})\bigr)$. 于是有一个满射 $W^{\oplus n'}\twoheadrightarrow V|_{G}$. 这表明 $V|_{G}$ 完全可约，同构于 $W^{\oplus e}$，其中 $e\leqslant n'$. 但此时 Schur 引理告诉我们 $n'=e$，最终我们证明了 $\operatorname{ev}_{W,V}$ 是一个同构. 由此还可推出 $W'$ 关于 $G'$ 不可约，且 $V|_{G'}\simeq W'^{\oplus n}$，其中 $n=\dim_{\mathbb C}(W)$.
这样我们证明了 ii). 把这个结论用于 $W\otimes_{\mathbb C}W'$ 的一个非零不可约子表示，也可推出 i).
剩下要证明 iii). 但若 $W=\bigoplus_{i}W_{i}$ 且 $W'=\bigoplus_{j}W'_{j}$，其中诸 $W_{i}$ 与诸 $W'_{j}$ 不可约，那么 $W\otimes W'=\bigoplus_{i,j}W_{i}\otimes W'_{j}$ 便是一个由不可约表示组成的直和分解.
\end{proof}
\end{numpara}

\begin{corollary}\label{cor:lie-1-12-10}
设 $\mathfrak g$ 是一个 $\mathbb C$-Lie 代数.
\begin{enumerate}[label=\roman*)]
\item 若 $\mathfrak g$ 的``全纯''表示完全可约，那么它的所有复表示都完全可约.
\item $\mathfrak g$ 的不可约表示具有形状 $(V,\rho)\otimes(V',\overline{\rho}')$，其中 $(V,\rho)$ 与 $(V',\rho')$ 是两个全纯不可约表示，而 $\overline{\rho}'$ 是 $\rho'$ 的共轭表示（因而它是``\newterm{lie-antiholomorphic}{反全纯的}{antiholomorphic}''，即 $\mathbb C$-半线性的）.
\end{enumerate}
\end{corollary}
\begin{proof}
想法是利用通过表示的限制给出的范畴等价 $\operatorname{Rep}^{\mathrm{hol}}_{\mathbb C}(\mathfrak g_{\mathbb C})\xrightarrow{\ \sim\ }\operatorname{Rep}_{\mathbb C}(\mathfrak g)$. 用 $\overline{\mathfrak g}$ 记这个 $\mathbb C$-Lie 代数 $\mathfrak g$，但其上的 $\mathbb C$-作用是 $\mathfrak g$ 上作用的共轭. 于是``恒等''映射是 $\mathfrak g$ 与 $\overline{\mathfrak g}$ 之间的一个半线性同构，而 $\overline{\mathfrak g}$ 的``全纯''表示恰是 $\mathfrak g$ 的``反全纯''表示. 现在，态射
\[
\begin{array}{rcl}
\mathfrak g_{\mathbb C}=\mathbb C\otimes_{\mathbb R}\mathfrak g & \longrightarrow & \mathfrak g\oplus\overline{\mathfrak g}\\
z\otimes X & \longmapsto & (zX,\overline{z}X)
\end{array}
\]
是 $\mathbb C$-Lie 代数之间的同构. 根据前面的命题，$\mathfrak g_{\mathbb C}$ 的任何``全纯''不可约表示都具有形状 $V\otimes_{\mathbb C}\overline{V}$，其中 $V$ 是 $\mathfrak g$ 的全纯不可约表示，而 $\overline{V}$ 是 $\overline{\mathfrak g}$ 的全纯不可约表示. 限制到 $\mathfrak g\subset\mathfrak g_{\mathbb C}$ 上便得 ii). i) 可类似证明.
\end{proof}
